\documentclass{article}
\usepackage{amsmath,amssymb}
\usepackage{xurl}
\usepackage[hidelinks]{hyperref}
% Shared commands for notes in docs/paper-gaps/.

\DeclareUnicodeCharacter{2080}{\ifmmode{}_{0}\else\textsubscript{0}\fi}
\DeclareUnicodeCharacter{2081}{\ifmmode{}_{1}\else\textsubscript{1}\fi}
\DeclareUnicodeCharacter{2082}{\ifmmode{}_{2}\else\textsubscript{2}\fi}
\DeclareUnicodeCharacter{2099}{\ifmmode{}_{n}\else\textsubscript{n}\fi}

\newcommand{\Lean}{\textsc{Lean}}
\ExplSyntaxOn
\NewDocumentCommand{\leanid}{m}{
  \ifmmode
    \textsf{\detokenize{#1}}
  \else
    \group_begin:
    \urlstyle{sf}
    \tl_set:Nn \l_tmpa_tl {#1}
    \tl_replace_all:Nnn \l_tmpa_tl {\_} {_}
    \exp_args:NV \path \l_tmpa_tl
    \group_end:
  \fi
}
\ExplSyntaxOff
\providecommand{\tr}{\operatorname{tr}}
\newcommand{\tnleanIssueBase}{https://github.com/LionSR/TNLean/issues/}
\newcommand{\tnleanPrBase}{https://github.com/LionSR/TNLean/pull/}
\newcommand{\tnleanDocBase}{https://sirui-lu.com/TNLean/docs/}
\newcommand{\ghissue}[1]{\href{\tnleanIssueBase#1}{GitHub issue \##1}}
\newcommand{\ghpr}[1]{\href{\tnleanPrBase#1}{PR \##1}}
\newcommand{\doclink}[1]{\href{\tnleanDocBase#1.html}{\path{#1}}}

% Dirac notation and the identity operator, as in blueprint/src/macros/common.tex.
% \providecommand keeps note-local definitions authoritative.
\usepackage{dsfont}
\providecommand{\ket}[1]{|#1\rangle}
\providecommand{\bra}[1]{\langle #1|}
\providecommand{\braket}[2]{\langle #1 | #2 \rangle}
\providecommand{\ketbra}[2]{|#1\rangle\!\langle #2|}
\providecommand{\Id}{\mathds{1}}
\providecommand{\one}{\mathds{1}}
\providecommand{\C}{\mathbb{C}}
\providecommand{\MN}[1]{M_{#1}(\C)}
\providecommand{\MD}{M_D(\C)}

% Verdict marker: \gapnote{<kind>}{<status>}. Kind is the policy
% classification (clarification, local-correction, scope-restriction,
% unfaithful, false-source, open-gap); status is open, wip (elimination
% underway), resolved, or historical. Machine-read by the site index;
% prints a verdict line.
\newcommand{\gapnote}[2]{\par\noindent\textbf{Verdict:} #1 (#2).\par}

\providecommand{\gapnote}[2]{\par\noindent\textbf{Verdict:} #1 (#2).\par}
\title{Coordinate Slices in Hastings' Higher-Dimensional Hierarchy}
\author{}
\date{October 28, 2026}
\begin{document}
\maketitle
\gapnote{scope-restriction}{resolved}

\section{The source and the coordinate setting}

Hastings, \emph{Solving gapped Hamiltonians locally}
\cite{Hastings2006Solving}, arXiv:cond-mat/0508554v3, source
\path{proj3new.tex}, lines 745--756, gives two possible extensions of the
hierarchical construction to higher dimensions: group sites according to
one coordinate, or combine $2^D$ sites at a time. It asserts that a
projection length polynomial in $\log N$ still suffices. The first option
is now constructed on full finite rectangular boxes in every dimension.

The more general construction uses a finite nonempty site set $V$, a
coordinate $z:V\to\mathbb N$, a radius $\ell\ge1$, and physical operators
$M_i$ supported on $\{v:|z(v)-z(i)|\le\ell\}$. Metric-ball locality gives
this coordinate support whenever the coordinate map does not increase
distances. On a rectangular box with the coordinate maximum metric, that
property follows from coordinate projection and is proved directly.
Metrics dominating the coordinate maximum metric also satisfy it.

\section{The pair operators are constructed}

Put $b(v)=\lfloor z(v)/(2\ell)\rfloor$ and
$b_i=\lfloor\max\{0,z(i)-\ell\}/(2\ell)\rfloor$. Every supporting site
of $M_i$ belongs to slice $b_i$ or $b_i+1$. For $B\ge1$ coarse slices,
assign the term to edge $b_i+1$ when $b_i+1<B$, and to edge zero
otherwise. Let $A_j$ be the sum of terms assigned to $j$. Then
\[
 \sum_{j<B}A_j=\sum_iM_i,
\]
and each internal $A_j$ is supported on slices $j-1,j$. If
$z(v)<2\ell B$ for every site, edge zero is supported on the last and
first slices. The latter support statement and coverage of all physical
sites are both proved, rather than inferred from the names of the
blocks.\footnote{Declarations:
\leanid{TNLean.MPS.AreaLaw.LocalSolution.coordinateSliceOperator},
\leanid{TNLean.MPS.AreaLaw.LocalSolution.sum_coordinateSliceOperator},
\leanid{TNLean.MPS.AreaLaw.LocalSolution.coordinateSliceOperator_mem_supportedOperators},
\leanid{TNLean.MPS.AreaLaw.LocalSolution.coordinateSliceOperator_mem_supportedOperators_of_metric}, and
\leanid{TNLean.MPS.AreaLaw.LocalSolution.coordinateSliceOperator_closing_mem_supportedOperators}.}

Write $n=|V|$. For a unit reference state $\psi$, suppose the Hermitian
local terms satisfy
\[
 \|M_i\|\le J,\qquad
 \|(M_i-\langle\psi,M_i\psi\rangle)\psi\|\le Je^{-r}.
\]
This is the local residual input of equation (low), lines 633--639, with
a common coefficient for the norm and residual estimates. Fixed local
constants are included in that coefficient. Summation over a fibre gives
\[
 \|A_j-\langle\psi,A_j\psi\rangle\|\le2nJ,
 \qquad
 \|(A_j-\langle\psi,A_j\psi\rangle)\psi\|\le nJe^{-r}.
\]
Thus no residual or support estimate for unspecified pair operators is
an assumption of this construction.\footnote{Declarations:
\leanid{TNLean.MPS.AreaLaw.LocalSolution.residual_edgePartition_le},
\leanid{TNLean.MPS.AreaLaw.LocalSolution.centeredNorm_edgePartition_le_two_mul},
\leanid{TNLean.MPS.AreaLaw.LocalSolution.residual_coordinateSliceOperator_le}, and
\leanid{TNLean.MPS.AreaLaw.LocalSolution.centeredNorm_coordinateSliceOperator_le_two_mul}.}

\section{The retained state and the physical gap}

Set $K=J_{\rm eff}=2nJ$, let
$h=\max\{2,\lceil\log_2 B\rceil\}$, and extend the finite edge family by
zero up to $2^h$ edges. The corrected hierarchy and its closing window
are determined by these actual operators. With
\[
 r\ge512h^2\bigl(1+\max\{0,\log(2nJ/\Delta)\}\bigr),
\]
they retain a unit vector $\phi$ with
\[
 \|\phi-\psi\|\le2(3^{h+1}-1)e^{-r/8}.
\]
The vector is retained at every node of the binary tree, and the finite
slice family covers all physical sites. The deficiency correction and
the extra closing window are as in the notes
\href{hastings05_hierarchical_projection.pdf}{on hierarchical projection}
and \href{hastings05_hierarchical_energy.pdf}{on hierarchical energy}.

If the physical Hamiltonian $H$ has ground projection $P_0$, ground
energy $E_0$, and gap $\Delta>0$, if $\psi$ is its ground state, and if
$\|H-\sum_iM_i\|\le nJe^{-r}$, then every orthonormal family in the
constructed final space has at most $2\operatorname{rank}P_0$ members.
The gap is that of $H$ itself.\footnote{Declarations:
\leanid{TNLean.MPS.AreaLaw.LocalSolution.coordinateSliceProjection},
\leanid{TNLean.MPS.AreaLaw.LocalSolution.exists_unit_mem_coordinateSliceProjection}, and
\leanid{TNLean.MPS.AreaLaw.LocalSolution.card_le_two_mul_re_trace_of_coordinateSliceProjection},
in \doclink{TNLean/MPS/AreaLaw/LocalSolution/CoordinateSliceApproximation}.}

\section{A polynomial logarithmic length on full boxes}

For $V=\prod_{k=1}^D\{0,\ldots,L_k-1\}$, with $D\ge1$ and all
$L_k\ge1$, choose any axis $a$ and take $B=L_a$. Varying that coordinate
while fixing the others at zero injects its side into $V$, so $B\le n$.
The coordinate extent and the metric-to-strip implication are immediate
from the box coordinates and are proved formally. The sufficient
projection length is therefore at most
\[
 512(\lceil\log_2n\rceil+2)^3
 \bigl(1+\max\{0,\log(2J/\Delta)\}\bigr).
\]
This proves the asserted polynomial logarithmic dependence for the
coordinate-slice option on full rectangular boxes in arbitrary dimension.
No bond-dimension formula is asserted by the source passage or supplied
by this estimate.\footnote{Declarations:
\leanid{TNLean.MPS.AreaLaw.LocalSolution.coordinateSide_le_card_box},
\leanid{TNLean.MPS.AreaLaw.LocalSolution.coordinateBox_slice_extent},
\leanid{TNLean.MPS.AreaLaw.LocalSolution.coordinateBox_ball_subset_strip}, and
\leanid{TNLean.MPS.AreaLaw.LocalSolution.hierarchyProjectionLength_coordinateBox_le},
\leanid{TNLean.MPS.AreaLaw.LocalSolution.exists_unit_mem_coordinateSliceProjection_of_box}, and
\leanid{TNLean.MPS.AreaLaw.LocalSolution.card_le_two_mul_re_trace_of_coordinateSliceProjection_of_box},
in \doclink{TNLean/MPS/AreaLaw/LocalSolution/CoordinateSliceBox}.}

The scope is explicit: the full-box corollary does not assert a bound on
the coordinate extent of arbitrary sparse embeddings, nor identify the
coordinate metric with the intrinsic distance on every lattice. The
abstract construction remains valid whenever its stated coordinate
support, coverage, and coarse-count bounds are established. The other
source option, grouping $2^D$ sites at each stage, is not used.

\bibliographystyle{unsrt}
\bibliography{references}
\end{document}
